\section{Search in a Binary Search Tree}
\label{sec:bst-search}

\problemstatement{Given the root of a binary search tree and a target
value $X$, find and return the subtree rooted at the node whose value
equals $X$, or return null if no such node exists.}

\begin{patternbox}
\emph{``Search for a value in a BST''} is the tree-shaped restatement of
ordinary binary search (the Binary Search chapter's opening section): the BST invariant
plays exactly the role sortedness played there, letting a single
comparison at each node discard an entire subtree (half the remaining
candidates) rather than requiring us to check it at all.
\end{patternbox}

\subsection*{Understanding the problem carefully}

At any node, comparing $X$ to that node's value tells us everything we
need: if $X$ equals the node's value, we are done. If $X$ is smaller, the
BST invariant guarantees that \emph{every} value in the right subtree is
even larger than the current node, so $X$ cannot possibly be there --
discard the right subtree entirely and recurse left. Symmetrically, if
$X$ is larger, discard the left subtree and recurse right.

\begin{ideabox}[One Comparison, One Subtree Discarded]
At each node: if the node is null, $X$ is not present (return null). If
$X$ equals the node's value, return this node. If $X<$ the node's value,
recurse into the left child. If $X>$ the node's value, recurse into the
right child.
\end{ideabox}

\subsection*{Why this is exactly binary search, with depth playing the role of $\log n$ steps}

Each step travels one level deeper into the tree and provably eliminates
an entire subtree from consideration, using the same discard argument as
array-based binary search: the eliminated subtree's values are
\emph{all} on the wrong side of $X$, guaranteed by the BST invariant, so
none of them need to be checked individually. Because each step descends
exactly one level, the total number of steps is bounded by the tree's
height -- $O(\log n)$ for a balanced tree, though note this degrades to
$O(n)$ for a degenerate (linked-list-shaped) tree, a distinction worth
remembering throughout this chapter.

\subsection*{Algorithm}

\begin{enumerate}
  \item If the current node is null: return null.
  \item If $X =$ current node's value: return the current node.
  \item If $X <$ current node's value: recurse into the left child.
  \item Else: recurse into the right child.
\end{enumerate}

\begin{algorithm}[H]
\caption{Search in a BST}
\begin{algorithmic}[1]
\Procedure{SearchBST}{\textit{root}, $X$}
  \If{$\textit{root} = \text{null}$} \Return \text{null} \EndIf
  \If{$X = \textit{root}.\textit{val}$} \Return \textit{root} \EndIf
  \If{$X < \textit{root}.\textit{val}$}
    \State \Return \Call{SearchBST}{$\textit{root}.\textit{left}$, $X$}
  \Else
    \State \Return \Call{SearchBST}{$\textit{root}.\textit{right}$, $X$}
  \EndIf
\EndProcedure
\end{algorithmic}
\end{algorithm}

\subsection*{Dry run}

\dryrun{Take the BST
$\begin{smallmatrix} & & 5 & & \\ & 3 & & 8 & \\ 1 & & 4 & & 9
\end{smallmatrix}$ (root $5$; left child $3$ with children $1,4$; right
child $8$ with right child $9$), searching for $X=4$.}

\begin{itemize}
  \item At $5$: $4<5$, go left to $3$.
  \item At $3$: $4>3$, go right to $4$.
  \item At $4$: $4=4$, found. Return this node.
\end{itemize}

\subsection*{C++ implementation}

\begin{lstlisting}
#include <bits/stdc++.h>
using namespace std;

struct TreeNode {
    int val;
    TreeNode* left;
    TreeNode* right;
    TreeNode(int v) : val(v), left(nullptr), right(nullptr) {}
};

TreeNode* searchBST(TreeNode* root, int X) {
    if (root == nullptr || root->val == X) return root;
    if (X < root->val) return searchBST(root->left, X);
    return searchBST(root->right, X);
}

int main() {
    TreeNode* root = new TreeNode(5);
    root->left = new TreeNode(3);
    root->right = new TreeNode(8);
    root->left->left = new TreeNode(1);
    root->left->right = new TreeNode(4);
    root->right->right = new TreeNode(9);

    TreeNode* found = searchBST(root, 4);
    cout << (found ? found->val : -1) << endl;
    return 0;
}
\end{lstlisting}

\begin{complexitybox}
\textbf{Time Complexity.} $O(h)$, where $h$ is the tree's height --
$O(\log n)$ for a balanced BST, $O(n)$ in the worst case for a
degenerate, linked-list-shaped tree.
\textbf{Space Complexity.} $O(h)$ for the recursion stack (or $O(1)$ if
rewritten iteratively, since no backtracking is ever needed).
\end{complexitybox}

\begin{takeawaybox}
Every BST operation in this chapter that walks from the root toward a
single target (search, insert, floor/ceil, predecessor/successor)
inherits this exact $O(h)$ complexity and this exact ``one comparison,
one subtree discarded'' justification. Internalizing it once here saves
re-deriving it in every later section.
\end{takeawaybox}
